\section[Алгоритмы на основе произведения матриц для всех пар вершин]{Алгоритмы на основе произведения матриц для всех пар вершин}
	\tikzsetfigurename{CFPQ_MatrixBased_}
\sectionmark{Алгоритмы на основе произведения матриц для всех пар вершин}
	\tikzsetfigurename{CFPQ_MatrixBased_}
\label{sec:CFPQ_MatrixBased}

В разделе~\ref{sec:CFPQ_Hellings}~был изложен алгоритм для решения задачи КС достижимости на основе CYK.
Заметим, что обход матрицы напоминает умножение матриц в ячейках которых множества нетерминалов:
\begin{align*}
M_3 = &M_1 \times M_2 \\
M_3[i,j] = &\sum_{k=1}^{n} M[i,k] * M[k,j]
\end{align*}
, где сумма~--- это объединение множеств:
\[
\sum_{k=1}^{n} = \bigcup_{k=1}^{n}
\]
, а поэлементное умножение определено следующим образом:
\[
S_1 * S_2 = \{N_1^0 ... N_1^m\} * \{N_2^0 ... N_2^l\} = \{N_3 \mid (N_3 \rightarrow N_1^i N_2^j) \in P\}.
\]
Таким образом, алгоритм решения задачи КС достижимости может быть выражена в терминах перемножения матриц над полукольцом с соответствующими операциями.

Для частного случая этой задачи, синтаксического анализа линейного входа, существует алгоритм Валианта~\cite{Valiant:1975:GCR:1739932.1740048}, использующий эту идею.
Однако он не обобщается на графы из-за того, что существенно использует возможность упорядочить обход матрицы (см. разницу в CYK для линейного случая и для графа).
Поэтому, хотя для линейного случая алгоритм Валианта является алгоритмом синтаксического анализа для произвольных КС грамматик за субкубическое время, его обобщение до задачи КС достижимости в произвольных графах с сохранением асимптотики является нетривиальной задачей~\cite{Yannakakis}.
В настоящее время алгоритм с субкубической сложностью получен только для частного случая~--- языка Дика с одним типом скобок~--- Филипом Брэдфордом~\cite{8249039}.

В случае с линейным входом, отдельного внимания заслуживает работа Лиллиан Ли (Lillian Lee)~\cite{Lee:2002:FCG:505241.505242}, где она показывает, что задача перемножения матриц сводима к синтаксическому анализу линейного входа. Аналогичных результатов для графов на текущий момент не известно.

Поэтому рассмотрим более простую идею, изложенную в статье Рустама Азимова~\cite{Azimov:2018:CPQ:3210259.3210264}: будем строить транзитивное замыкание графа через наивное (не через возведение в квадрат) умножение матриц.

Пусть $\mathcal{G} = (V, E)$~--- входной граф и $G = (N,\Sigma,P)$~--- входная грамматика. Тогда алгоритм может быть сформулирован, как представлено в листинге~\ref{alg:graphParse}.

\begin{algorithm}[H]
%TODO!!!
%\begin{algorithmic}[1]
%\caption{Context-free recognizer for graphs}
%\label{alg:graphParse}
%\Function{contextFreePathQuerying}{$\mathcal{G}$, G}

%    \State{$n \gets$ количество узлов в $\mathcal{G}$}
%    \State{$E \gets$ направленные ребра в $\mathcal{G}$}
%    \State{$P \gets$ набор продукций из $G$}
%    \State{$T \gets$ матрица $n \times n$, в которой каждый элемент $\varnothing$}
%    \ForAll{$(i,x,j) \in E$}
%    \Comment{Инициализация матрицы}
%        \State{$T_{i,j} \gets T_{i,j} \cup \{A~|~(A \rightarrow x) \in P \}$}
%    \EndFor
%    \ForAll{$i \in 0\ldots n-1$}
%    \Comment{Добавление петель для нетерминалов, порождающих пустую строку}
%        \State{$T_{i,i} \gets T_{i,i} \cup \{ A \in N \mid A \to \varepsilon \}$}
%    \EndFor
%    \While{матрица $T$ меняется}

%        \State{$T \gets T \cup (T \times T)$}
%        \Comment{Вычисление транзитивного замыкания}
%    \EndWhile
%\State \Return $T$
%\EndFunction
%\end{algorithmic}
\end{algorithm}


\begin{example}[Пример работы]

Пусть есть граф $\mathcal{G}$:
\begin{center}
  \input{figures/edge_labelled_graph_a_3_loop_b_2_loop}
\end{center}

и грамматика $G$:
\begin{align*}
S   &\to A B    &A  \to a \\
S  &\to A S_1   &B  \to b\\
S_1 &\to S B
\end{align*}


Пусть $T_i$~--- матрица, полученная из $T$ после применения цикла, описанного в строках \textbf{8-9} алгоритма~\ref{alg:graphParse}, $i$ раз.
Тогда $T_0$, полученная напрямую из графа, выглядит следующим образом:

\[
T_0 = \begin{pmatrix}
    \varnothing & \{A\}       & \varnothing & \{B\}       \\
    \varnothing & \varnothing & \{A\}       & \varnothing \\
    \{A\}       & \varnothing & \varnothing & \varnothing \\
    \{B\}       & \varnothing & \varnothing & \varnothing \\
\end{pmatrix}
\]

Далее показано получение матрицы $T_1$.

\[
T_0 \times T_0 = \begin{pmatrix}
    \varnothing & \varnothing & \varnothing & \varnothing \\
    \varnothing & \varnothing & \varnothing & \varnothing \\
    \varnothing & \varnothing & \varnothing & \{S\}       \\
    \varnothing & \varnothing & \varnothing & \varnothing \\
\end{pmatrix}
\]

\[
T_1 = T_0 \cup (T_0 \times T_0) = \begin{pmatrix}
    \varnothing & \{A\}       & \varnothing & \{B\}       \\
    \varnothing & \varnothing & \{A\}       & \varnothing \\
    \{A\}       & \varnothing & \varnothing & \cellcolor{lightgray} \{\pmb{S}\}       \\
    \{B\}       & \varnothing & \varnothing & \varnothing \\
\end{pmatrix}
\]

После первой итерации цикла нетерминал в ячейку $T[2,3]$ добавился нетерминал $S$.
Это означает, что существует такой путь $\pi$ из вершины 2 в вершину 3 в графе $\mathcal{G}$, что $S \xrightarrow{*} \omega(\pi)$. В данном примере путь состоит из двух ребер $2 \xrightarrow{a} 0$ и $ 0 \xrightarrow{b} 3$, так что $S \xrightarrow{*} ab$.

Вычисление транзитивного замыкания заканчивается через $k$ итераций, когда достигается неподвижная точка процесса: $T_{k-1} = T_k$. Для данного примера $k = 13$, так как $T_{13} = T_{12}$. Весь процесс работы алгоритма (все матрицы $T_i$) показан ниже (на каждой итерации новые элементы выделены жирным).

{\footnotesize
\begin{alignat*}{7}
& &&T_2 &&= \begin{pmatrix}
\varnothing & \{A\}       & \varnothing & \{B\}       \\
\varnothing & \varnothing & \{A\}       & \varnothing \\
\cellcolor{lightgray} \{A, \pmb{S_1}\}  & \varnothing & \varnothing & \{S\}       \\
\{B\}       & \varnothing & \varnothing & \varnothing \\
\end{pmatrix} \ \ \ \ &&T_3 &&= \begin{pmatrix}
\varnothing & \{A\}       & \varnothing & \{B\}       \\
\cellcolor{lightgray} \{\pmb{S}\}       & \varnothing & \{A\}       & \varnothing \\
\{A, S_1\}  & \varnothing & \varnothing & \{S\}       \\
\{B\}       & \varnothing & \varnothing & \varnothing \\
\end{pmatrix} \\ & &&T_4 &&= \begin{pmatrix}
\varnothing & \{A\}       & \varnothing & \{B\}       \\
\{S\}       & \varnothing & \{A\}       & \cellcolor{lightgray} \{\pmb{S_1}\}     \\
\{A, S_1\}  & \varnothing & \varnothing & \{S\}       \\
\{B\}       & \varnothing & \varnothing & \varnothing \\
\end{pmatrix}  \ \ \ \ &&T_5 &&= \begin{pmatrix}
\varnothing & \{A\}       & \varnothing & \cellcolor{lightgray} \{B, \pmb{S}\}    \\
\{S\}       & \varnothing & \{A\}       & \{S_1\}     \\
\{A, S_1\}  & \varnothing & \varnothing & \{S\}       \\
\{B\}       & \varnothing & \varnothing & \varnothing \\
\end{pmatrix} \\ & &&T_6 &&= \begin{pmatrix}
\cellcolor{lightgray} \{\pmb{S_1}\}     & \{A\}       & \varnothing & \{B, S\}    \\
\{S\}       & \varnothing & \{A\}       & \{S_1\}     \\
\{A, S_1\}  & \varnothing & \varnothing & \{S\}       \\
\{B\}       & \varnothing & \varnothing & \varnothing \\
\end{pmatrix} \ \ \ \ &&T_7 &&= \begin{pmatrix}
\{S_1\}     & \{A\}       & \varnothing & \{B, S\}    \\
\{S\}       & \varnothing & \{A\}       & \{S_1\}     \\
\cellcolor{lightgray} \{A, S_1, \pmb{S}\}  & \varnothing & \varnothing & \{S\}    \\
\{B\}       & \varnothing & \varnothing & \varnothing \\
\end{pmatrix}  \\
& &&T_8 &&= \begin{pmatrix}
\{S_1\}     & \{A\}       & \varnothing & \{B, S\}    \\
\{S\}       & \varnothing & \{A\}       & \{S_1\}     \\
\{A, S_1, S\}  & \varnothing & \varnothing & \cellcolor{lightgray} \{S, \pmb{S_1}\} \\
\{B\}       & \varnothing & \varnothing & \varnothing \\
\end{pmatrix} \ \ \ \ &&T_9 &&= \begin{pmatrix}
\{S_1\}     & \{A\}       & \varnothing & \{B, S\}    \\
\{S\}       & \varnothing & \{A\}       & \cellcolor{lightgray} \{S_1, \pmb{S}\}     \\
\{A, S_1, S\}  & \varnothing & \varnothing & \{S, S_1\} \\
\{B\}       & \varnothing & \varnothing & \varnothing \\
\end{pmatrix} \\ & &&T_{10} &&= \begin{pmatrix}
\{S_1\}     & \{A\}       & \varnothing & \{B, S\}    \\
\cellcolor{lightgray} \{S, \pmb{S_1}\}       & \varnothing & \{A\}       & \{S_1, S\}     \\
\{A, S_1, S\}  & \varnothing & \varnothing & \{S, S_1\} \\
\{B\}       & \varnothing & \varnothing & \varnothing \\
\end{pmatrix}  \ \ \ \  &&T_{11} &&= \begin{pmatrix}
\cellcolor{lightgray} \{S_1, \pmb{S}\}     & \{A\}       & \varnothing & \{B, S\}    \\
\{S, S_1\}       & \varnothing & \{A\}       & \{S_1, S\}     \\
\{A, S_1, S\}  & \varnothing & \varnothing & \{S, S_1\} \\
\{B\}       & \varnothing & \varnothing & \varnothing \\
\end{pmatrix} \\ & &&T_{12} &&= \begin{pmatrix}
\{S_1, S\}     & \{A\}       & \varnothing & \cellcolor{lightgray} \{B, S, \pmb{S_1}\}    \\
\{S, S_1\}       & \varnothing & \{A\}       & \{S_1, S\}     \\
\{A, S_1, S\}  & \varnothing & \varnothing & \{S, S_1\} \\
\{B\}       & \varnothing & \varnothing & \varnothing \\
\end{pmatrix} \ \ \ \ &&T_{13} &&= \begin{pmatrix}
\{S_1, S\}     & \{A\}       & \varnothing & \{B, S, S_1\}    \\
\{S, S_1\}       & \varnothing & \{A\}       & \{S_1, S\}     \\
\{A, S_1, S\}  & \varnothing & \varnothing & \{S, S_1\} \\
\{B\}       & \varnothing & \varnothing & \varnothing \\
\end{pmatrix}
\end{alignat*}
}

Таким образом, результат алгоритма~\ref{alg:graphParse} для нашего примера~--- это матрица $T_{13} = T_{12}$. Заметим, что для данного алгоритма приведённый пример также является худшим случаем: на каждой итерации в матрицу добавляется ровно один нетерминал, при том, что необходимо заполнить порядка $O(n^2)$ ячеек.

\end{example}


%% FIXME: Исправить подраздел
% \subsection{Расширение алгоритма для конъюнктивных грамматик}

% Матричный алгоритм для конъюнктивных грамматик отличается от алгоритма~\ref{alg:graphParse} для контекстно-свободных грамматик только операцией умножения матриц, в остальном алгоритм остается без изменений. Определим операцию умножения матриц.
% \begin{definition}
%     Пусть $M_1$ и $M_2$ матрицы размера $n$. Определим операцию $\circ$  следующим образом:
%      \[M_1 \circ M_2 = M_3,\]
%      \[M_3 [i,j] = \{A \mid \exists (A \rightarrow B_1 C_1 \& \ldots \& B_m C_m) \in P, (B_k , C_k) \in d[i,j] \forall k = 1,\ldots,m\}\],
%      где \[d[i,j] = \bigcup_{k = 1}^{n} M_1 [i,k] \times M_2 [k,j].\]
% \end{definition}

% Важно заметить, что алгоритм для конъюнктивных грамматик, в отличие от алгоритма для контекстно-свободных грамматик, дает лишь верхнюю оценку ответа. То есть все нетерминалы, которые должны быть в ячейках матрицы результата, содержатся там, но вместе с ними содержатся и лишние нетерминалы. Рассмотрим пример, иллюстрирующий появление лишних нетерминалов.

% \begin{example}
%     Грамматика $G$:
%     \begin{align*}
%     S &\to AB \& DC & C &\to c \\
%     A &\to a        & D &\to DC \mid b\\
%     B &\to BC \mid b
%     \end{align*}
%     Очевидно, что грамматика $G$ задает язык из одного слова $L(G) = \{abc\} = \{abc^*\} \cap \{a^* bc\}$.

%     Пусть есть граф $\mathcal{G}$:
%     \begin{center}
%       \input{figures/multi/graph0.tex}
%     \end{center}
%     Применяя алгоритм, получим, что существует путь из вершины 0 в вершину 4, выводимый из нетерминала $S$. Однако очевидно, что в графе такого пути нет.
%     Такое поведение алгоритма наблюдается из-за того, что существует путь ``abcc'', соответствующий $L(AB) = \{abc^*\}$ и путь ``aabc'', соответствующий $L(DC) = \{a^{*}bc\}$, но они различны. Однако алгоритм не может это проверить, так как оперирует понятием достижимости между вершинами, а не наличием различных путей. Более того, в общем случае для конъюнктивных грамматик такую проверку реализовать нельзя. Поэтому для классической семантики достижимости с ограничениями в терминах конъюнктивных грамматик результат работы алгоритма является оценкой сверху.

%     Существует альтернативная семантика, когда мы трактуем конъюнкцию в правой части правил как конъюнкцию в Datalog (подробнее о Datalog в параграфе~\ref{Subsection Datalog}). Т.е если есть правило $S \to AB \& DC$, то должен быть путь принадлежащий языку $L(AB)$ и путь принадлежащий языку $L(DC)$. В такой семантике алгоритм дает точный ответ.
% \end{example}

% Подробнее алгоритм описан в статье Рустама Азимова и Семёна Григорьева~\cite{565CECD7E8F5C6063935B41DB41797AA37D53B04}. Стоит также отметить, что обобщения данного алгоритма для булевых грамматик не существует, хотя и существует частное решение для случая, когда граф не содержит циклов (является DAG-ом), предложенное Екатериной Шеметовой~\cite{Shemetova2019}.
\mytodo{Особенности реализации перенесены в раздел~\ref{sec:Conclusion_ImplFeatures} Заключения.}
